\documentclass{article}
\usepackage{amsmath}
\usepackage{geometry}
\geometry{paperheight=29.7cm, paperwidth=21cm, margin=1cm}

\begin{document}

\section*{Q1.}
What will be the ratio of distance moved by a freely falling body from rest in $4^{th}$ and $5^{th}$ second of journey.
    
\section*{Q2.}
A wooden block of dimensions 20 cm $\times$ 10 cm $\times$ 5 cm has a mass of 300 g. Find the thrust and pressure exerted by the block on a horizontal surface when it is placed on the surface with its sides of (a) 20 cm $\times$ 10 cm, (b) 20 cm $\times$ 5 cm, and (c) 10 cm $\times$ 5 cm.

\section*{Q3.}
A metal sphere of radius 5 cm and mass 2 kg is completely immersed in water. Find the buoyant force acting on the sphere and the apparent weight of the sphere in water. Take the density of water as $1000 \, \text{kg/m}^3$ and $g = 10 \, \text{m/s}^2$.

\section*{Q4.}
A piece of iron weighs 200 g in air and 180 g in water. Find the relative density of iron and the volume of the piece of iron. Take the density of water as $1000 \, \text{kg/m}^3$.

\section*{Q5.}
A solid cube of side 10 cm and density 800 kg/m\(^3\) floats in a liquid of density 1000 kg/m\(^3\). Find the height of the cube above the liquid surface and the fraction of the volume of the cube submerged in the liquid.

\section*{Q6.}
What is the escape velocity of a rocket launched from the surface of the earth? Assume that the radius of the earth is 6400 km and the acceleration due to gravity is $9.8 \, \text{m/s}^2$.

\section*{Q7.}
What is the orbital period of a satellite that orbits the earth at a height of 200 km above the surface of the earth? Assume that the radius of the earth is 6400 km and the acceleration due to gravity is $9.8 \, \text{m/s}^2$.

\section*{Q8.}
Kepler's third law states that the square of the time it takes a planet to orbit the sun is related to the cube of its average distance from the sun. For Earth, with an orbital period (time to orbit the sun) of 365 days and an average distance of 150 million km, we can use the formula:

\[ T^2 \propto R^3 \]

Now, let's talk about Mars. If Mars is 1.5 times farther from the sun than Earth, what would be its orbital period? 

1. Use the formula \( T_{\text{mars}}^2 = T_{\text{earth}}^2 \times \left(\frac{R_{\text{mars}}}{R_{\text{earth}}}\right)^3 \) to calculate the orbital period of Mars (\(T_{\text{mars}}\)).
  
2. Given:
   - \( T_{\text{earth}} = 365 \) days
   - \( R_{\text{earth}} = 150 \) million km

3. Substitute these values into the formula and solve for \( T_{\text{mars}} \).


\end{document}
